% --- Archon physics formalization source begin ---
% archon:physics
% archon:covers PhyXMiniProblems/problem_phyx_mini_0743.lean
% archon:source-report reports/phyx_mini/problem_phyx_mini_0743.source.json

\chapter{Physics problem phyx\_mini\_0743}
\label{ch:PhyXMiniProblems_problem_phyx_mini_0743}

\paragraph{Problem source.}
\#\# Physical scenario

In 1939 or 1940, Emanuel Zacchini took his humancannonball act to an extreme: After being shot from a cannon, he soared over three Ferris wheels and into a net. Assume that he is launched with a speed of \$26.5\textbackslash{} m/s\$ and at an angle of \$53.0\textasciicircum{}\textbackslash{}circ\$.

\#\# Question

How far from the cannon should the net's center have been positioned (neglect air drag)?

\#\# Answer choices

- A: A: 63 m

- B: B: 66 m

- C: C: 69 m

- D: D: 72 m

\#\# Figure caption (auxiliary; use the image as primary evidence)

The image is a diagram illustrating a projectile motion scenario. Here's a detailed breakdown:

1. **Objects and Features:**
   - **Ramp/Launch Point:** On the left side, there's a sloped ramp leading to the launch point, with a wheel or circular object at the base.
   - **Ferris Wheels:** There are three Ferris wheels positioned along the trajectory path. They are of equal height and all appear to be 18 meters tall.
   - **Net:** On the right side, there is a net for catching the projectile, positioned 3 meters above the ground.

2. **Measurements and Text:**
   - The launch angle is labeled as \textbackslash{}( \textbackslash{}theta\_0 \textbackslash{}) and the initial velocity as \textbackslash{}( v\_0 \textbackslash{}).
   - The height from the ground to the launch point is 3 meters.
   - The distance from the launch point to the net is labeled as \textbackslash{}( R \textbackslash{}), with a known distance of 23 meters from the launch point to the first Ferris wheel tower.
   - The net is 3 meters above the ground where the projectile is aimed to land.

3. **Trajectory:**
   - The trajectory of the projectile is shown as a curved arrow that arcs over the Ferris wheels and towards the net.

This diagram is likely illustrating a physics problem involving kinematics and projectile motion.

\#\# Dataset metadata

Kinematics; Physical Model Grounding Reasoning, Spatial Relation Reasoning

\paragraph{Recorded answer/context.}
C: C: 69 m

\paragraph{Figure/image path.}
/root/proposal\_for\_physic/science-mango/phyx\_mini\_run/phyx\_data/test\_image/743.png

\paragraph{Formalization target.}
create a compiling Lean file with sorry bodies at `PhyXMiniProblems/problem\_phyx\_mini\_0743.lean`. The Lean declarations must preserve the physical quantities, dimensions or dimensional roles, figure labels, governing-law hypotheses, and final relation expressed by this problem.
Use Mathlib/Physlib names found through LeanExplore where available. If a physics API is missing, introduce faithful local abstractions rather than scalar placeholder aliases.

\begin{theorem}[Physics formalization target]
\label{thm:physics:phyx_mini_0743:target}
\lean{PhyXMiniProblems.ProblemPhyXMini0743.problem_phyx_mini_0743}
\uses{def:physics:phyx-mini-0743:phyxminiproblems-problemphyxmini0743-lengthinmeters, def:physics:phyx-mini-0743:phyxminiproblems-problemphyxmini0743-humancannonballsetup, def:physics:phyx-mini-0743:phyxminiproblems-problemphyxmini0743-matchesprimaryfigure, def:physics:phyx-mini-0743:phyxminiproblems-problemphyxmini0743-matchesproblemreadouts, def:physics:phyx-mini-0743:phyxminiproblems-problemphyxmini0743-usesstandardnearearthgravity, def:physics:phyx-mini-0743:phyxminiproblems-problemphyxmini0743-hasphysicalcannonballparameters, def:physics:phyx-mini-0743:phyxminiproblems-problemphyxmini0743-satisfiesidealprojectilemotion, def:physics:phyx-mini-0743:phyxminiproblems-problemphyxmini0743-clearsallthreeferriswheels, def:physics:phyx-mini-0743:phyxminiproblems-problemphyxmini0743-computedrangeinmeters, def:physics:phyx-mini-0743:phyxminiproblems-problemphyxmini0743-matchesdisplayeddistance, def:physics:phyx-mini-0743:phyxminiproblems-problemphyxmini0743-isuniqueclosestdisplayeddistance}
The assigned autoformalize agent should translate this physics problem into Lean declarations in the covered file, with theorem and lemma proof bodies written as `by sorry`.
\end{theorem}
\begin{proof}
This is an autoformalization task, not a proof task. Produce statements that can later be proved by the physics prover without weakening the physical model.
\end{proof}

% --- Archon declaration topology begin ---
\section{Declaration topology}
\begin{definition}[acceleration Dimension]
  \label{def:physics:phyx-mini-0743:phyxminiproblems-problemphyxmini0743-accelerationdimension}
  \lean{PhyXMiniProblems.ProblemPhyXMini0743.accelerationDimension}
  The physical dimension of acceleration, length * time⁻²
\end{definition}
\begin{proof}
  This is the declared model definition of acceleration Dimension.
\end{proof}
\begin{definition}[Length Quantity]
  \label{def:physics:phyx-mini-0743:phyxminiproblems-problemphyxmini0743-lengthquantity}
  \lean{PhyXMiniProblems.ProblemPhyXMini0743.LengthQuantity}
  A nonnegative, unit-independent physical length
\end{definition}
\begin{proof}
  This is the declared model definition of Length Quantity.
\end{proof}
\begin{definition}[Signed Length Quantity]
  \label{def:physics:phyx-mini-0743:phyxminiproblems-problemphyxmini0743-signedlengthquantity}
  \lean{PhyXMiniProblems.ProblemPhyXMini0743.SignedLengthQuantity}
  A signed one-dimensional physical coordinate or displacement
\end{definition}
\begin{proof}
  This is the declared model definition of Signed Length Quantity.
\end{proof}
\begin{definition}[Time Quantity]
  \label{def:physics:phyx-mini-0743:phyxminiproblems-problemphyxmini0743-timequantity}
  \lean{PhyXMiniProblems.ProblemPhyXMini0743.TimeQuantity}
  A nonnegative, unit-independent elapsed time
\end{definition}
\begin{proof}
  This is the declared model definition of Time Quantity.
\end{proof}
\begin{definition}[Acceleration Quantity]
  \label{def:physics:phyx-mini-0743:phyxminiproblems-problemphyxmini0743-accelerationquantity}
  \lean{PhyXMiniProblems.ProblemPhyXMini0743.AccelerationQuantity}
  \uses{def:physics:phyx-mini-0743:phyxminiproblems-problemphyxmini0743-accelerationdimension}
  A nonnegative, unit-independent acceleration magnitude
\end{definition}
\begin{proof}
  This is the declared model definition of Acceleration Quantity.
\end{proof}
\begin{definition}[nonnegative SIReadout]
  \label{def:physics:phyx-mini-0743:phyxminiproblems-problemphyxmini0743-nonnegativesireadout}
  \lean{PhyXMiniProblems.ProblemPhyXMini0743.nonnegativeSIReadout}
  Read a nonnegative dimensionful quantity in coherent SI units
\end{definition}
\begin{proof}
  This is the declared model definition of nonnegative SIReadout.
\end{proof}
\begin{definition}[length In Meters]
  \label{def:physics:phyx-mini-0743:phyxminiproblems-problemphyxmini0743-lengthinmeters}
  \lean{PhyXMiniProblems.ProblemPhyXMini0743.lengthInMeters}
  \uses{def:physics:phyx-mini-0743:phyxminiproblems-problemphyxmini0743-lengthquantity, def:physics:phyx-mini-0743:phyxminiproblems-problemphyxmini0743-nonnegativesireadout}
  Metre readout of a physical length
\end{definition}
\begin{proof}
  This is the declared model definition of length In Meters.
\end{proof}
\begin{definition}[signed Length In Meters]
  \label{def:physics:phyx-mini-0743:phyxminiproblems-problemphyxmini0743-signedlengthinmeters}
  \lean{PhyXMiniProblems.ProblemPhyXMini0743.signedLengthInMeters}
  \uses{def:physics:phyx-mini-0743:phyxminiproblems-problemphyxmini0743-signedlengthquantity}
  Metre readout of a signed coordinate
\end{definition}
\begin{proof}
  This is the declared model definition of signed Length In Meters.
\end{proof}
\begin{definition}[time In Seconds]
  \label{def:physics:phyx-mini-0743:phyxminiproblems-problemphyxmini0743-timeinseconds}
  \lean{PhyXMiniProblems.ProblemPhyXMini0743.timeInSeconds}
  \uses{def:physics:phyx-mini-0743:phyxminiproblems-problemphyxmini0743-timequantity, def:physics:phyx-mini-0743:phyxminiproblems-problemphyxmini0743-nonnegativesireadout}
  Second readout of a physical elapsed time
\end{definition}
\begin{proof}
  This is the declared model definition of time In Seconds.
\end{proof}
\begin{definition}[speed In Meters Per Second]
  \label{def:physics:phyx-mini-0743:phyxminiproblems-problemphyxmini0743-speedinmeterspersecond}
  \lean{PhyXMiniProblems.ProblemPhyXMini0743.speedInMetersPerSecond}
  Metre-per-second readout of a physical speed
\end{definition}
\begin{proof}
  This is the declared model definition of speed In Meters Per Second.
\end{proof}
\begin{definition}[acceleration In Meters Per Second Squared]
  \label{def:physics:phyx-mini-0743:phyxminiproblems-problemphyxmini0743-accelerationinmeterspersecondsquared}
  \lean{PhyXMiniProblems.ProblemPhyXMini0743.accelerationInMetersPerSecondSquared}
  \uses{def:physics:phyx-mini-0743:phyxminiproblems-problemphyxmini0743-accelerationquantity, def:physics:phyx-mini-0743:phyxminiproblems-problemphyxmini0743-nonnegativesireadout}
  Metre-per-second-squared readout of an acceleration magnitude
\end{definition}
\begin{proof}
  This is the declared model definition of acceleration In Meters Per Second Squared.
\end{proof}
\begin{definition}[degrees]
  \label{def:physics:phyx-mini-0743:phyxminiproblems-problemphyxmini0743-degrees}
  \lean{PhyXMiniProblems.ProblemPhyXMini0743.degrees}
  Convert a degree readout into a physical angle modulo 2π
\end{definition}
\begin{proof}
  This is the declared model definition of degrees.
\end{proof}
\begin{definition}[Projectile Model]
  \label{def:physics:phyx-mini-0743:phyxminiproblems-problemphyxmini0743-projectilemodel}
  \lean{PhyXMiniProblems.ProblemPhyXMini0743.ProjectileModel}
  The idealization named in the prose
\end{definition}
\begin{proof}
  This is the declared model definition of Projectile Model.
\end{proof}
\begin{definition}[Ferris Wheel]
  \label{def:physics:phyx-mini-0743:phyxminiproblems-problemphyxmini0743-ferriswheel}
  \lean{PhyXMiniProblems.ProblemPhyXMini0743.FerrisWheel}
  The three distinct Ferris wheels crossed from left to right
\end{definition}
\begin{proof}
  This is the declared model definition of Ferris Wheel.
\end{proof}
\begin{definition}[Figure Object]
  \label{def:physics:phyx-mini-0743:phyxminiproblems-problemphyxmini0743-figureobject}
  \lean{PhyXMiniProblems.ProblemPhyXMini0743.FigureObject}
  \uses{def:physics:phyx-mini-0743:phyxminiproblems-problemphyxmini0743-ferriswheel}
  Named physical objects visible in image 743.png
\end{definition}
\begin{proof}
  This is the declared model definition of Figure Object.
\end{proof}
\begin{definition}[Figure Label]
  \label{def:physics:phyx-mini-0743:phyxminiproblems-problemphyxmini0743-figurelabel}
  \lean{PhyXMiniProblems.ProblemPhyXMini0743.FigureLabel}
  Numerical and symbolic labels printed in the primary image
\end{definition}
\begin{proof}
  This is the declared model definition of Figure Label.
\end{proof}
\begin{definition}[Figure Anchor]
  \label{def:physics:phyx-mini-0743:phyxminiproblems-problemphyxmini0743-figureanchor}
  \lean{PhyXMiniProblems.ProblemPhyXMini0743.FigureAnchor}
  \uses{def:physics:phyx-mini-0743:phyxminiproblems-problemphyxmini0743-ferriswheel}
  Schematic points needed to retain the spatial relations in the image
\end{definition}
\begin{proof}
  This is the declared model definition of Figure Anchor.
\end{proof}
\begin{definition}[Human Cannonball Figure]
  \label{def:physics:phyx-mini-0743:phyxminiproblems-problemphyxmini0743-humancannonballfigure}
  \lean{PhyXMiniProblems.ProblemPhyXMini0743.HumanCannonballFigure}
  \uses{def:physics:phyx-mini-0743:phyxminiproblems-problemphyxmini0743-figureobject, def:physics:phyx-mini-0743:phyxminiproblems-problemphyxmini0743-figurelabel, def:physics:phyx-mini-0743:phyxminiproblems-problemphyxmini0743-figureanchor}
  Qualitative and schematic-coordinate evidence from the primary raster. Coordinates are dimensionless drawing coordinates, not measured lengths
\end{definition}
\begin{proof}
  This is the declared model definition of Human Cannonball Figure.
\end{proof}
\begin{definition}[Planar Position]
  \label{def:physics:phyx-mini-0743:phyxminiproblems-problemphyxmini0743-planarposition}
  \lean{PhyXMiniProblems.ProblemPhyXMini0743.PlanarPosition}
  \uses{def:physics:phyx-mini-0743:phyxminiproblems-problemphyxmini0743-signedlengthquantity}
  A physical point in the vertical plane of the flight
\end{definition}
\begin{proof}
  This is the declared model definition of Planar Position.
\end{proof}
\begin{definition}[Human Cannonball Setup]
  \label{def:physics:phyx-mini-0743:phyxminiproblems-problemphyxmini0743-humancannonballsetup}
  \lean{PhyXMiniProblems.ProblemPhyXMini0743.HumanCannonballSetup}
  \uses{def:physics:phyx-mini-0743:phyxminiproblems-problemphyxmini0743-lengthquantity, def:physics:phyx-mini-0743:phyxminiproblems-problemphyxmini0743-timequantity, def:physics:phyx-mini-0743:phyxminiproblems-problemphyxmini0743-accelerationquantity, def:physics:phyx-mini-0743:phyxminiproblems-problemphyxmini0743-projectilemodel, def:physics:phyx-mini-0743:phyxminiproblems-problemphyxmini0743-ferriswheel, def:physics:phyx-mini-0743:phyxminiproblems-problemphyxmini0743-humancannonballfigure, def:physics:phyx-mini-0743:phyxminiproblems-problemphyxmini0743-planarposition}
  Independent physical quantities and observables for the cannonball flight. The horizontal net distance labelled R is deliberately an unconstrained physical field; the governing laws below relate it to the trajectory
\end{definition}
\begin{proof}
  This is the declared model definition of Human Cannonball Setup.
\end{proof}
\begin{definition}[Matches Primary Figure]
  \label{def:physics:phyx-mini-0743:phyxminiproblems-problemphyxmini0743-matchesprimaryfigure}
  \lean{PhyXMiniProblems.ProblemPhyXMini0743.MatchesPrimaryFigure}
  \uses{def:physics:phyx-mini-0743:phyxminiproblems-problemphyxmini0743-humancannonballsetup}
  Primary-image evidence. It retains all named objects and labels, the equal schematic heights of launch and net center, vertical wheel towers, and the left-to-right ordering. It contains no numerical readout for R
\end{definition}
\begin{proof}
  This is the declared model definition of Matches Primary Figure.
\end{proof}
\begin{definition}[Matches Problem Readouts]
  \label{def:physics:phyx-mini-0743:phyxminiproblems-problemphyxmini0743-matchesproblemreadouts}
  \lean{PhyXMiniProblems.ProblemPhyXMini0743.MatchesProblemReadouts}
  \uses{def:physics:phyx-mini-0743:phyxminiproblems-problemphyxmini0743-lengthinmeters, def:physics:phyx-mini-0743:phyxminiproblems-problemphyxmini0743-speedinmeterspersecond, def:physics:phyx-mini-0743:phyxminiproblems-problemphyxmini0743-degrees, def:physics:phyx-mini-0743:phyxminiproblems-problemphyxmini0743-humancannonballsetup}
  Numerical quantities stated in the prose or printed in the figure
\end{definition}
\begin{proof}
  This is the declared model definition of Matches Problem Readouts.
\end{proof}
\begin{definition}[Uses Standard Near Earth Gravity]
  \label{def:physics:phyx-mini-0743:phyxminiproblems-problemphyxmini0743-usesstandardnearearthgravity}
  \lean{PhyXMiniProblems.ProblemPhyXMini0743.UsesStandardNearEarthGravity}
  \uses{def:physics:phyx-mini-0743:phyxminiproblems-problemphyxmini0743-accelerationinmeterspersecondsquared, def:physics:phyx-mini-0743:phyxminiproblems-problemphyxmini0743-humancannonballsetup}
  Standard near-Earth gravitational acceleration used by the answer key
\end{definition}
\begin{proof}
  This is the declared model definition of Uses Standard Near Earth Gravity.
\end{proof}
\begin{definition}[Has Physical Cannonball Parameters]
  \label{def:physics:phyx-mini-0743:phyxminiproblems-problemphyxmini0743-hasphysicalcannonballparameters}
  \lean{PhyXMiniProblems.ProblemPhyXMini0743.HasPhysicalCannonballParameters}
  \uses{def:physics:phyx-mini-0743:phyxminiproblems-problemphyxmini0743-lengthinmeters, def:physics:phyx-mini-0743:phyxminiproblems-problemphyxmini0743-timeinseconds, def:physics:phyx-mini-0743:phyxminiproblems-problemphyxmini0743-speedinmeterspersecond, def:physics:phyx-mini-0743:phyxminiproblems-problemphyxmini0743-accelerationinmeterspersecondsquared, def:physics:phyx-mini-0743:phyxminiproblems-problemphyxmini0743-humancannonballsetup}
  Positivity, the upward/rightward launch branch, and the physical ordering of the three wheels and net. These conditions constrain geometry but do not assign the requested range a numerical value
\end{definition}
\begin{proof}
  This is the declared model definition of Has Physical Cannonball Parameters.
\end{proof}
\begin{definition}[Satisfies Ideal Projectile Motion]
  \label{def:physics:phyx-mini-0743:phyxminiproblems-problemphyxmini0743-satisfiesidealprojectilemotion}
  \lean{PhyXMiniProblems.ProblemPhyXMini0743.SatisfiesIdealProjectileMotion}
  \uses{def:physics:phyx-mini-0743:phyxminiproblems-problemphyxmini0743-timequantity, def:physics:phyx-mini-0743:phyxminiproblems-problemphyxmini0743-lengthinmeters, def:physics:phyx-mini-0743:phyxminiproblems-problemphyxmini0743-signedlengthinmeters, def:physics:phyx-mini-0743:phyxminiproblems-problemphyxmini0743-timeinseconds, def:physics:phyx-mini-0743:phyxminiproblems-problemphyxmini0743-speedinmeterspersecond, def:physics:phyx-mini-0743:phyxminiproblems-problemphyxmini0743-accelerationinmeterspersecondsquared, def:physics:phyx-mini-0743:phyxminiproblems-problemphyxmini0743-humancannonballsetup}
  The no-drag, constant-gravity projectile equations in coherent SI readouts. The final two fields identify the positive flight-time trajectory point with the net center. No field contains 69, an answer label, or the eliminated level-flight range formula
\end{definition}
\begin{proof}
  This is the declared model definition of Satisfies Ideal Projectile Motion.
\end{proof}
\begin{definition}[Clears All Three Ferris Wheels]
  \label{def:physics:phyx-mini-0743:phyxminiproblems-problemphyxmini0743-clearsallthreeferriswheels}
  \lean{PhyXMiniProblems.ProblemPhyXMini0743.ClearsAllThreeFerrisWheels}
  \uses{def:physics:phyx-mini-0743:phyxminiproblems-problemphyxmini0743-timequantity, def:physics:phyx-mini-0743:phyxminiproblems-problemphyxmini0743-lengthinmeters, def:physics:phyx-mini-0743:phyxminiproblems-problemphyxmini0743-signedlengthinmeters, def:physics:phyx-mini-0743:phyxminiproblems-problemphyxmini0743-timeinseconds, def:physics:phyx-mini-0743:phyxminiproblems-problemphyxmini0743-humancannonballsetup}
  The prose says Zacchini passes above all three wheels. This observation is kept as clearance data at genuine positive trajectory times; it neither fixes the net distance nor asserts the requested answer
\end{definition}
\begin{proof}
  This is the declared model definition of Clears All Three Ferris Wheels.
\end{proof}
\begin{lemma}[net Range eq level Flight Formula]
  \label{lem:physics:phyx-mini-0743:phyxminiproblems-problemphyxmini0743-netrange-eq-levelflightformula}
  \lean{PhyXMiniProblems.ProblemPhyXMini0743.netRange_eq_levelFlightFormula}
  \uses{def:physics:phyx-mini-0743:phyxminiproblems-problemphyxmini0743-lengthinmeters, def:physics:phyx-mini-0743:phyxminiproblems-problemphyxmini0743-speedinmeterspersecond, def:physics:phyx-mini-0743:phyxminiproblems-problemphyxmini0743-accelerationinmeterspersecondsquared, def:physics:phyx-mini-0743:phyxminiproblems-problemphyxmini0743-humancannonballsetup, def:physics:phyx-mini-0743:phyxminiproblems-problemphyxmini0743-matchesproblemreadouts, def:physics:phyx-mini-0743:phyxminiproblems-problemphyxmini0743-hasphysicalcannonballparameters, def:physics:phyx-mini-0743:phyxminiproblems-problemphyxmini0743-satisfiesidealprojectilemotion}
  Eliminating the nonzero flight time when the launch and target heights agree gives the standard level-flight range. This is a derived conclusion, not a premise of the numerical answer theorem
\end{lemma}
\begin{proof}
  The conclusion follows from the governing relations and figure readouts named in the statement of net Range eq level Flight Formula.
\end{proof}
\begin{definition}[computed Range In Meters]
  \label{def:physics:phyx-mini-0743:phyxminiproblems-problemphyxmini0743-computedrangeinmeters}
  \lean{PhyXMiniProblems.ProblemPhyXMini0743.computedRangeInMeters}
  \uses{def:physics:phyx-mini-0743:phyxminiproblems-problemphyxmini0743-degrees}
  Exact ideal-model range from the stated speed, angle, and gravity
\end{definition}
\begin{proof}
  This is the declared model definition of computed Range In Meters.
\end{proof}
\begin{definition}[Answer Choice]
  \label{def:physics:phyx-mini-0743:phyxminiproblems-problemphyxmini0743-answerchoice}
  \lean{PhyXMiniProblems.ProblemPhyXMini0743.AnswerChoice}
  The four answer labels printed with the problem
\end{definition}
\begin{proof}
  This is the declared model definition of Answer Choice.
\end{proof}
\begin{definition}[distance In Meters]
  \label{def:physics:phyx-mini-0743:phyxminiproblems-problemphyxmini0743-answerchoice-distanceinmeters}
  \lean{PhyXMiniProblems.ProblemPhyXMini0743.AnswerChoice.distanceInMeters}
  \uses{def:physics:phyx-mini-0743:phyxminiproblems-problemphyxmini0743-answerchoice}
  Whole-metre distance printed beside each answer label
\end{definition}
\begin{proof}
  This is the declared model definition of distance In Meters.
\end{proof}
\begin{definition}[recorded Answer Choice]
  \label{def:physics:phyx-mini-0743:phyxminiproblems-problemphyxmini0743-recordedanswerchoice}
  \lean{PhyXMiniProblems.ProblemPhyXMini0743.recordedAnswerChoice}
  \uses{def:physics:phyx-mini-0743:phyxminiproblems-problemphyxmini0743-answerchoice}
  The answer label recorded in the supplied dataset metadata
\end{definition}
\begin{proof}
  This is the declared model definition of recorded Answer Choice.
\end{proof}
\begin{definition}[Matches Displayed Distance]
  \label{def:physics:phyx-mini-0743:phyxminiproblems-problemphyxmini0743-matchesdisplayeddistance}
  \lean{PhyXMiniProblems.ProblemPhyXMini0743.MatchesDisplayedDistance}
  \uses{def:physics:phyx-mini-0743:phyxminiproblems-problemphyxmini0743-lengthinmeters, def:physics:phyx-mini-0743:phyxminiproblems-problemphyxmini0743-humancannonballsetup, def:physics:phyx-mini-0743:phyxminiproblems-problemphyxmini0743-answerchoice}
  Agreement with a displayed distance rounded to the nearest metre
\end{definition}
\begin{proof}
  This is the declared model definition of Matches Displayed Distance.
\end{proof}
\begin{definition}[Is Closest Displayed Distance]
  \label{def:physics:phyx-mini-0743:phyxminiproblems-problemphyxmini0743-isclosestdisplayeddistance}
  \lean{PhyXMiniProblems.ProblemPhyXMini0743.IsClosestDisplayedDistance}
  \uses{def:physics:phyx-mini-0743:phyxminiproblems-problemphyxmini0743-lengthinmeters, def:physics:phyx-mini-0743:phyxminiproblems-problemphyxmini0743-humancannonballsetup, def:physics:phyx-mini-0743:phyxminiproblems-problemphyxmini0743-answerchoice}
  A displayed choice is at least as close as every other printed value
\end{definition}
\begin{proof}
  This is the declared model definition of Is Closest Displayed Distance.
\end{proof}
\begin{definition}[Is Unique Closest Displayed Distance]
  \label{def:physics:phyx-mini-0743:phyxminiproblems-problemphyxmini0743-isuniqueclosestdisplayeddistance}
  \lean{PhyXMiniProblems.ProblemPhyXMini0743.IsUniqueClosestDisplayedDistance}
  \uses{def:physics:phyx-mini-0743:phyxminiproblems-problemphyxmini0743-humancannonballsetup, def:physics:phyx-mini-0743:phyxminiproblems-problemphyxmini0743-answerchoice, def:physics:phyx-mini-0743:phyxminiproblems-problemphyxmini0743-isclosestdisplayeddistance}
  A displayed choice is the unique closest printed value
\end{definition}
\begin{proof}
  This is the declared model definition of Is Unique Closest Displayed Distance.
\end{proof}
% --- Archon declaration topology end ---
% --- Archon physics formalization source end ---
